\section{Implementation}

\subsection{Data Sources}

\subsubsection{Census Data}

Population data comes from the 2020 Census Redistricting Data (PL 94-171) Summary File \cite{census2020pl}. This provides total population counts at the census tract level—the same data states use for official redistricting. We use 84,414 census tracts covering all 50 states.

Geographic boundaries come from Census TIGER/Line shapefiles (2020 vintage) \cite{tiger-census}, which define tract polygons with coordinate precision suitable for adjacency determination and boundary length computation.

\subsubsection{Enacted Districts}

For baseline comparison, we use the Census Bureau's TIGER/Line shapefiles for the 118th Congressional Districts \cite{census2023tigerline}. These represent enacted 2020 Census-based boundaries used in the 2022 and 2024 elections. This dataset provides official boundaries for rigorous compactness comparison.

\subsubsection{Political Data}

For political analysis (Section 6), we use 2020 presidential election results at the census tract level from MIT Election Data + Science Lab \cite{mit-electionlab}. We compute partisan lean as Biden vote share: $\text{Lean}_d = \text{Biden votes} / (\text{Biden} + \text{Trump})$ for district $d$.

\subsection{Software Implementation}

\subsubsection{Core Algorithm}

The recursive bisection algorithm is implemented in Python 3.13+ using:
\begin{itemize}
\item \textbf{GeoPandas} for geographic data processing
\item \textbf{Shapely} for geometric operations (boundary intersections)
\item \textbf{NetworkX} for adjacency graph construction (verification only)
\item \textbf{METIS} (gpmetis.exe) for graph partitioning
\end{itemize}

The main partitioning logic resides in \texttt{src/apportionment/partition/recursive\_bisection.py}, which implements both baseline and edge-weighted modes via a \texttt{partition\_mode} parameter.

\subsubsection{Edge Weight Computation}

Edge weights are computed by \texttt{src/apportionment/data/adjacency.py}:

\begin{verbatim}
def compute_edge_weights(tracts_gdf):
    """Compute boundary lengths for adjacent tracts."""
    edge_weights = {}

    # Build spatial index for efficiency
    sindex = tracts_gdf.sindex

    for idx, tract in tracts_gdf.iterrows():
        # Find potential neighbors
        candidates = sindex.query(
            tract.geometry.bounds,
            predicate='intersects'
        )

        for neighbor_idx in candidates:
            if neighbor_idx <= idx:
                continue

            neighbor = tracts_gdf.loc[neighbor_idx]

            # Compute intersection
            intersection = tract.geometry.boundary \
                .intersection(neighbor.geometry.boundary)

            if intersection.is_empty:
                # Water crossing - use median
                weight = median_land_boundary
            elif intersection.geom_type == 'Point':
                # Corner contact
                weight = 0.1
            else:
                # Real shared boundary
                weight = intersection.length

            edge_weights[(idx, neighbor_idx)] = weight

    return edge_weights
\end{verbatim}

\subsubsection{METIS Integration}

METIS input files are generated in CSR format. For edge-weighted mode:

\begin{verbatim}
def write_metis_graph(graph, edge_weights, filename):
    """Write METIS format with edge weights."""
    n_vertices = len(graph.nodes)
    n_edges = len(graph.edges)

    with open(filename, 'w') as f:
        # Header: format 011 = edge weights + vertex weights
        f.write(f"{n_vertices} {n_edges} 011 1\n")

        for node in sorted(graph.nodes):
            # Vertex weight (population)
            vweight = graph.nodes[node]['population']
            f.write(f"{vweight}")

            # Neighbors and edge weights
            for neighbor in sorted(graph.neighbors(node)):
                eweight = edge_weights.get(
                    (min(node, neighbor),
                     max(node, neighbor)),
                    1
                )
                # Convert to centimeters
                eweight_cm = max(1, int(eweight * 100))
                # METIS uses 1-indexed nodes
                f.write(f" {neighbor + 1} {eweight_cm}")

            f.write("\n")
\end{verbatim}

METIS is invoked via subprocess with appropriate flags for each mode.

\subsection{Pipeline Architecture}

The complete pipeline consists of:

\begin{enumerate}
\item \textbf{Data loading}: Load census tracts with populations
\item \textbf{Adjacency computation}: Build graph of neighboring tracts
\item \textbf{Edge weight computation} (edge-weighted mode only): Calculate boundary lengths
\item \textbf{Recursive partitioning}: Partition into districts
\item \textbf{District assignment}: Map each tract to final district
\item \textbf{Geometry construction}: Dissolve tract boundaries to form district polygons
\item \textbf{Compactness calculation}: Compute Polsby-Popper and Reock scores
\item \textbf{Political analysis}: Aggregate election results by district
\item \textbf{Visualization}: Generate maps and summary statistics
\end{enumerate}

The pipeline processes all 50 states in parallel, with each state running independently. Total runtime for 50 states: $\sim$2-3 hours on a modern workstation.

\subsection{Computational Performance}

\begin{table}[h]
\centering
\caption{Performance by State Size}
\begin{tabular}{lrrr}
\toprule
\textbf{State} & \textbf{Tracts} & \textbf{Districts} & \textbf{Time} \\
\midrule
Wyoming & 141 & 1 & < 1 sec \\
New Hampshire & 269 & 2 & 8 sec \\
Iowa & 792 & 4 & 22 sec \\
Ohio & 3,493 & 15 & 89 sec \\
Texas & 5,834 & 38 & 184 sec \\
California & 9,213 & 52 & 267 sec \\
\bottomrule
\end{tabular}
\end{table}

Performance scales approximately linearly with number of tracts. Edge weight computation adds 10-30 seconds of preprocessing overhead but does not affect METIS runtime significantly.

Memory usage peaks at $\sim$2 GB for California (largest state). The limiting factor is loading high-resolution TIGER/Line geometries, not graph partitioning itself.

\subsection{Reproducibility}

All code and data are available in the project repository. To reproduce results:

\begin{enumerate}
\item Install dependencies: Python 3.13+, GeoPandas, METIS
\item Download 2020 Census PL 94-171 data and TIGER/Line shapefiles
\item Run pipeline: \texttt{python run\_redistricting.py --year 2020 --version v1 --partition-mode edge-weighted}
\item Compare against enacted: \texttt{python compare\_vs\_enacted.py}
\end{enumerate}

Fixed random seeds ensure deterministic results. Given identical inputs, the algorithm produces byte-identical outputs.

\subsection{Coordinate Reference Systems}

Accurate boundary length computation requires projected coordinate systems with equal-area properties. We use:

\begin{itemize}
\item \textbf{Contiguous 48 states}: Albers Equal Area (EPSG:5070)
\item \textbf{Alaska}: Alaska Albers (EPSG:3338)
\item \textbf{Hawaii}: Hawaii NAD83 (EPSG:2784)
\end{itemize}

These projections minimize distortion for their respective regions, providing meter-accurate boundary length measurements. Computing edge weights in geographic coordinates (NAD83 lat/lon) would yield incorrect results due to degree-based measurements.
